\section{Planejamento}

O desenvolvimento desta pesquisa está estruturado nas seguintes etapas, com atividades distribuídas ao longo dos próximos meses, conforme cronograma previsto:

\begin{itemize}
    \item \textbf{Etapa 1: Preparação de Dados}
    \begin{itemize}
        \item \textbf{Junho}: Coleta, pré-processamento e integração dos dados espaciais (redes, POIs, condicionantes ambientais e socioeconômicos).
    \end{itemize}
    
    \item \textbf{Etapa 2: Implementação Computacional}
    \begin{itemize}
        \item \textbf{Julho}: Desenvolvimento dos scripts para aplicação das funções de penalização e cálculo dos indicadores de acessibilidade;
        \item \textbf{Julho}: Parametrização e testes dos algoritmos de rota e penalização.
    \end{itemize}
    
    \item \textbf{Etapa 3: Aplicação Experimental}
    \begin{itemize}
        \item \textbf{Agosto}: Execução dos experimentos nos estudos de caso (Curitiba e Porto Alegre);
        \item \textbf{Agosto}: Geração de mapas e síntese preliminar dos resultados de acessibilidade.
    \end{itemize}

    \item \textbf{Etapa 4: Análise Socioespacial e Síntese de Resultados}
    \begin{itemize}
        \item \textbf{Setembro}: Aplicação das análises de desigualdade, regressão e clusterização;
        \item \textbf{Setembro}: Interpretação e consolidação dos resultados.
    \end{itemize}
    
    \item \textbf{Etapa 5: Redação e Revisão}
    \begin{itemize}
        \item \textbf{Outubro}: Consolidação da coletânea de artigos para defesa;
    \end{itemize}
\end{itemize}

O cronograma prevê a conclusão de todas as etapas em um horizonte de cinco meses, com possíveis ajustes em função do andamento das etapas experimentais e análises complementares.